\originalchapter{干涉、衍射与薄膜}\label{ch:interference}
\originalsection{相位与相干}
两束光在同一点叠加的是电场矢量, 亮度则与时间平均场平方有关. 若同偏振、同频率、相位差为$\phi$且每束幅度$E_0$, 则
\[
E_1+E_2=2E_0\cos(\phi/2)\sin(\omega t+\phi/2),\quad
I=4I_1\cos^2(\phi/2).
\]
相位差为$2\pi m$时增强, 为$(2m+1)\pi$时相消. 若相位在观测时间内随机变化, 干涉交叉项平均为零, 只见两强度相加. 稳定可见条纹需要足够相干及匹配的偏振; 严格单色是简化条件, 实际有限谱宽在足够短光程差内也可干涉.

额外路程$\delta$产生相位差$2\pi\delta/\lambda$, 还需加初始相位及反射附加相位. 两窄缝间距$d$, 同相照明, 远场方向$\theta$的光程差近似为$d\sin\theta$. 因而亮纹满足$d\sin\theta=m\lambda$, 暗纹满足$d\sin\theta=(m+1/2)\lambda$. 屏距$D\gg d$, 若$|y|\ll D$, 则$\sin\theta\simeq y/D$, 亮纹位置$y_m\simeq m\lambda D/d$, 条纹间距$\lambda D/d$. 大角度时不能把$\tan\theta$和$\sin\theta$混用.

\originalsection{单缝与有限宽双缝}
衍射积分中的相位从几何得到. 设观察点距缝中心$r_0$, 横向缝坐标为$x'$, 相对孔径法线的方向角为$\theta$, 则
\[
r(x')=\sqrt{r_0^2-2r_0x'\sin\theta+x'^2}\simeq r_0-x'\sin\theta.
\]
幅度因子$1/r$可近似共同常数, 传播相位则每增加路程$\Delta r$改变$k\Delta r$. 以$e^{-i\omega t}$为时间约定, 删去共同相位后的小元相量可取$e^{-ikx'\sin\theta}$; 对对称缝换为其共轭得到同一强度, 下式采用正号. 二次路程误差的相位也须小, 足够的夫琅禾费条件是$r_0\gg a^2/\lambda$并且$r_0\gg a$. 对双缝用总孔径横向尺度$A_{\rm tot}$替代$a$, 要求$r_0\gg A_{\rm tot}^2/\lambda$. 这与屏上$|y|\ll D$的小角近似是不同条件.

欧拉公式把每个小元的正弦振幅与相位记成一个复数, 相量相加后取总振幅的模平方, 所以能用积分合成连续孔径. 这不是把各小元强度直接相加, 后者会删掉衍射的交叉项.

惠更斯--菲涅耳图像把波前各点看作次级波源, 再叠加其场. 对宽$a$的均匀照明单缝, 远场各点相位随横向坐标$x'$变化为$kx'\sin\theta$. 相量幅度相对于中央值为
\[
\frac1a\int_{-a/2}^{a/2}e^{ikx'\sin\theta}\dd x'
=\frac{\sin\beta}{\beta},\qquad \beta=\frac{\pi a\sin\theta}{\lambda}.
\]
为看清积分, 记$\eta=k\sin\theta$, 则$\int_{-a/2}^{a/2}e^{i\eta x^{\prime}}\dd x^{\prime}=[e^{i\eta x^{\prime}}/(i\eta)]_{-a/2}^{a/2}=2\sin(\eta a/2)/\eta$. 除以$a$后即是$\sin\beta/\beta$, $\eta=0$时直接积分为$a$. 因此$I(\theta)=I(0)(\sin\beta/\beta)^2$, 中央处取连续极限为1. 非中央暗纹满足$a\sin\theta=m\lambda$, $m$为非零整数. 中央亮区宽, 两侧旁瓣较弱. 这是均匀标量近似的远场推导, 不能用来代替任意近场传播或任意偏振边界问题.

两条同宽缝间距$d$, 相量还乘两缝间相位的合成因子, 得
\[
I(\theta)=I(0)\left(\frac{\sin\beta}{\beta}\right)^2
\cos^2\left(\frac{\pi d\sin\theta}{\lambda}\right).
\]
两缝中心在$\pm d/2$时, 各自单缝振幅再带$e^{\pm i\eta d/2}$, 相加为$2\cos(\eta d/2)$倍单缝幅度. 除以中央两缝总幅度后才出现归一化的$\cos^2$因子. $d$是缝中心距且$d>a$, 因而快速双缝条纹受较宽的单缝包络调制. 例如$d=2a$, 双缝第$\pm2$级亮纹正好落在单缝第$\pm1$暗纹处, 实际缺级, 不能只看双缝亮纹条件就断言该处非零. 相消点只是把能量重新分配到其他方向, 不说明总体能量消失. 若只取某一个观察点, $\langle(E_1+E_2)^2\rangle$一般不等于$\langle E_1^2\rangle+\langle E_2^2\rangle$, 能量守恒应对完整系统或截面分析.

\originalsection{薄膜与驻波}
肥皂泡和油膜的颜色来自膜上、下表面反射波的干涉. 折射率$n=c/v_{\rm phase}$描述相位传播速度, 膜内波长为$\lambda_0/n$, 因而法向往返厚度$2t$的传播相位为$4\pi nt/\lambda_0$. 再计反射的相位跳跃才决定增强或相消. 以下采用普通非吸收介质近法向入射的两主反射波近似. 空气中折射率$n>1$、厚$t$的膜, 近法向照射时膜内往返光程差为$2nt$. 从低折射率到高折射率反射电场有$\pi$相位反转, 从高到低没有. 若只有一次反转, 反射光增强条件为$2nt=(m+1/2)\lambda_0$, 相消条件为$2nt=m\lambda_0$, $\lambda_0$是真空波长. 若两界面都反转或都不反转, 两条件交换. 在反射相位规则适用的斜入射范围, 比较同一输出波前的传播相位对应光程差$2nt\cos\theta_t$, $\theta_t$为膜内折射角. 在两外侧介质均为空气时, Snell关系$\sin\theta_{\rm out}=n\sin\theta_t$给出$2nt/\cos\theta_t-2t\tan\theta_t\sin\theta_{\rm out}=2nt\cos\theta_t$. 它包含两束出射光的侧向路径修正, 不是把膜内几何路程$2t/\cos\theta_t$直接乘$n$. 对任意大角或特定偏振越过布儒斯特角时, 需重新核反射系数的相位, 不无条件沿用上述反转口诀. 各反射波幅度可能不同, “相消”未必完全黑暗.

一列波与反向传播的同频波叠加形成驻波, 例如$\sin(kx-\omega t)+\sin(kx+\omega t)=2\sin kx\cos\omega t$. 相邻节点距离$\lambda/2$. 微波实验用反射板制造驻波来量波长, 接收天线转动可检验偏振; CD表面规则间距则类似光栅产生分色. 这些可观察结果都从相位叠加而来, 不能只凭亮度照片断定某一条具体波的场值.


电磁驻波不能处处用行波的$B=E/c$. 两列沿$\pm x$、同偏振的波若在$x=0$满足金属面切向电场为零, 合成可写$E_y=2E_0\sin kx\cos\omega t$, $B_z=-2(E_0/c)\cos kx\sin\omega t$. 电场节点与磁场节点相差$\lambda/4$, 平均向前能流为零, 场能却随时间交换. 微波接收器若主要耦合电场, 测到的是电场振幅的节点间距, 不是把磁场与电场的所有极小位置混在一起.
